\section{A geometric introduction to the quantum constructions}
\label{sec:quantumintro}

This part uses several constructions that share the word
\emph{quantum} but have different definitions. We begin with a
linear-algebra question: how does one compare two ways of combining
three vector spaces that carry symmetries? The answer supplies the
local weights used later to build topological invariants. The opening
example needs only tensor products and changes of basis; the
representation-theoretic terms are introduced as they arise.
Later subsections also use elementary complex analysis and the
geometry already discussed. No physical interpretation is needed
for the finite sums or special-function identities.

The organizing operation is cutting and gluing. A decomposition gives
intermediate labels; changing the decomposition gives a change-of-basis
matrix or integral kernel. The quantum constructions make these
operations exact. In suitable large-parameter limits their equations
recover geometric gluing, and their leading exponents recover geometric
actions. The simplest complete example below is a finite sum for the
figure-eight knot whose growth gives the volume of two regular ideal
tetrahedra.

\subsection{Six labels: recoupling and tetrahedral weights}
\label{sec:introrecoupling}

\paragraph{Why start with representations?}
A triangulated three-manifold has many descriptions: its tetrahedra
can be assembled in different orders, and the triangulation itself
can be changed. To build an invariant from local data, we need
weights with exact rules for comparing these descriptions.
Representation theory provides such rules. Combining representations
uses tensor products; decomposing them into simpler pieces introduces
intermediate labels; changing the order of decomposition produces a
matrix. The identities satisfied by these matrices will become the
identities needed for changing a triangulation. This is the reason
for the algebra in this subsection.

The Lie algebra $\mathfrak{sl}_2=\mathfrak{sl}_2(\CC)$ is the vector
space of traceless $2\times2$ complex matrices, with bracket
$[X,Y]=XY-YX$. A \emph{representation} on a vector space $V$ assigns
to each $X$ a linear operator $\rho(X)$ on $V$, linearly in $X$,
so that $\rho([X,Y])=[\rho(X),\rho(Y)]$. It is
\emph{irreducible} if no nonzero proper subspace is preserved by all
these operators. Thus an irreducible representation is a building
block for a symmetry action.

The finite-dimensional irreducibles are
\[
 V_j=\operatorname{Sym}^{2j}(\CC^2),\qquad
 j=0,\tfrac12,1,\tfrac32,\ldots,\qquad \dim V_j=2j+1.
\]
Here $\operatorname{Sym}^m(\CC^2)$ is the space of symmetric
$m$-fold tensors on $\CC^2$; its basis has one element for each
possible number $0,\ldots,m$ of copies of the second standard basis
vector. In particular, $V_0=\CC$ is the trivial representation and
$V_{1/2}=\CC^2$ is the usual matrix action.
The name \emph{spin} for $j$ comes from rotations: the compact group
$\mathrm{SU}(2)$ double-covers $\mathrm{SO}(3)$, and the complexification
of its Lie algebra is $\mathfrak{sl}_2$. Integer spins descend to
$\mathrm{SO}(3)$; nonintegral half-integer spins need the double cover.
We use this familiar family because its tensor-product rule is
especially explicit. At this stage $j$ is a representation label,
not an edge length.

\paragraph{Combining two building blocks.}
On a tensor product the infinitesimal symmetry acts on both factors:
\[
 X(v\otimes w)=(Xv)\otimes w+v\otimes(Xw).
\]
The Clebsch--Gordan rule says
\[
 V_a\otimes V_b
   \cong\bigoplus_{s=|a-b|}^{a+b}V_s,
 \qquad s\text{ increases in steps of }1.
\]
Each allowed $V_s$ occurs once. Equivalently, the triple $(a,b,s)$
is \emph{admissible} when it satisfies the triangle inequalities
and $a+b+s$ is an integer. The direct-sum sign means that every
vector has a unique sum of components in these invariant subspaces.

For example,
\[
 V_{1/2}\otimes V_{1/2}=V_0\oplus V_1,\qquad 2\cdot2=1+3.
\]
Choose an orthonormal basis $e_+,e_-$ of $\CC^2$. The one-dimensional
summand is spanned by
\[
 \frac{e_+\otimes e_- - e_-\otimes e_+}{\sqrt2},
\]
and the three-dimensional summand is the symmetric subspace, with
basis
\[
 e_+\otimes e_+,\qquad
 \frac{e_+\otimes e_-+e_-\otimes e_+}{\sqrt2},\qquad
 e_-\otimes e_-.
\]
The antisymmetric vector is fixed by $\mathrm{SL}_2$ because its
transformation factor is the determinant, which is one.
This gives a concrete meaning to a decomposition into irreducibles.

\paragraph{What the symbols $B$ and $F$ mean.}
Now combine three factors and specify a desired output type $V_d$.
The vector space
\[
 H_d=\operatorname{Hom}_{\mathfrak{sl}_2}
       (V_d,V_a\otimes V_b\otimes V_c)
\]
consists of linear maps that commute with the symmetry action.
Such a map is called an \emph{intertwiner}. The dimension of $H_d$
counts how many copies of $V_d$ occur in the triple tensor product.

Choose, once and for all, a normalized inclusion
$C_s^{ab}:V_s\longrightarrow V_a\otimes V_b$ for each admissible
triple. For example, $C_0^{1/2,1/2}(1)$ can be the antisymmetric
unit vector just displayed. There are two ways to compose these
inclusions:
\[
 \begin{aligned}
 B_s^{(ab)c}&=(C_s^{ab}\otimes\mathrm{id}_{V_c})\,C_d^{sc},\\
 B_t^{a(bc)}&=(\mathrm{id}_{V_a}\otimes C_t^{bc})\,C_d^{at}.
 \end{aligned}
\]
Both are maps from $V_d$ to the same triple tensor product.
The parentheses record which two factors were combined first.
The index $s$ ranges over labels for which $(a,b,s)$ and $(s,c,d)$
are admissible; $t$ ranges over labels for which $(b,c,t)$ and
$(a,t,d)$ are admissible. Each collection is a basis of $H_d$.

The letter $B$ therefore denotes a \emph{basis intertwiner}.
The numbers $F_{st}^{abc;d}$ are defined to be the coefficients
that express the first basis in the second:
\begin{equation}
 B_s^{(ab)c}=\sum_t F_{st}^{abc;d}B_t^{a(bc)}.
 \label{eq:introrecoupling}
\end{equation}
For fixed $a,b,c,d$, this is an ordinary finite change-of-basis
matrix, with row index $s$ and column index $t$. The sum runs
over the allowed intermediate types, not over vectors of $V_t$.
Normalization and phase choices for the inclusions matter:
changing them changes the entries of $F$. One may use the
$\mathrm{SU}(2)$-invariant Hermitian inner products to make the
bases orthonormal; $F$ is then unitary.

\paragraph{A complete $2\times2$ example.}
Take $a=b=c=d=\tfrac12$. Applying the two-factor rule twice gives
\[
 V_{1/2}^{\otimes3}\cong V_{3/2}\oplus V_{1/2}\oplus V_{1/2}.
\]
Thus $H_{1/2}$ has dimension two. In the first basis we distinguish
the copies by whether the first pair has type $s=0$ or $s=1$;
in the second we use the type $t=0$ or $t=1$ of the last pair.

The whole calculation can be seen in three tensor coordinates.
Write
\[
 u=e_+\otimes e_+\otimes e_-,\qquad
 v=e_+\otimes e_-\otimes e_+,\qquad
 w=e_-\otimes e_+\otimes e_+.
\]
An intertwiner from $V_{1/2}$ is determined by its value on $e_+$.
For the following phase choices those values are
\[
 \begin{aligned}
 B_0^{(ab)c}(e_+)&=\frac{v-w}{\sqrt2},
 &B_1^{(ab)c}(e_+)&=\frac{2u-v-w}{\sqrt6},\\
 B_0^{a(bc)}(e_+)&=\frac{u-v}{\sqrt2},
 &B_1^{a(bc)}(e_+)&=\frac{u+v-2w}{\sqrt6}.
 \end{aligned}
\]
To see why these vectors describe the two copies, the raising
matrix $\left(\begin{smallmatrix}0&1\\0&0\end{smallmatrix}\right)$
sends each of $u,v,w$ to $e_+\otimes e_+\otimes e_+$ under the
tensor-product action. It kills precisely the combinations whose
three coefficients sum to zero. These form a two-dimensional
space; applying its transpose, the lowering matrix, supplies the corresponding
images of $e_-$. The vectors above give two orthonormal bases
of this same space. Their subscripts $0,1$ record the
antisymmetric or symmetric behavior of the chosen pair.

Taking their inner products gives
\[
 \begin{pmatrix}B_0^{(ab)c}\\ B_1^{(ab)c}\end{pmatrix}
 =\begin{pmatrix}
 -\tfrac12&\tfrac{\sqrt3}{2}\\
 \tfrac{\sqrt3}{2}&\tfrac12
 \end{pmatrix}
 \begin{pmatrix}B_0^{a(bc)}\\ B_1^{a(bc)}\end{pmatrix}.
\]
For instance, the upper-left entry is
$\langle(v-w)/\sqrt2,(u-v)/\sqrt2\rangle=-1/2$.
The rows have norm one and are perpendicular, as a change between
orthonormal bases requires. Although tensor products are associative,
the bases selected by the two decompositions are different.
\emph{Recoupling} is the operation of converting between them.
There is no asymptotic approximation in this example.

\paragraph{From a basis change to a tetrahedron.}
The coefficients in \eqref{eq:introrecoupling} depend on six
labels $a,b,c,d,s,t$. With conventional dimension and sign factors
they are the \emph{$6j$-symbols}. The name records these six
representation labels; it does not refer to the size of the matrix.
The four admissible triples
\[
 (a,b,s),\quad(s,c,d),\quad(b,c,t),\quad(a,t,d)
\]
can be assigned to the four faces of a tetrahedron. Explicitly,
on a tetrahedron with vertices $1,2,3,4$, put
$a,b,s,d,c,t$ on edges $12,23,13,14,34,24$, respectively.
Every face then carries one of the four triples.
This explains the tetrahedral arrangement of a recoupling coefficient.
Its interpretation as a \emph{metric} tetrahedron requires an
additional asymptotic statement.

\paragraph{Deforming the rules.}
A quantum group deforms this tensor-product calculus. For
$U_{q_{\rm alg}}(\mathfrak{sl}_2)$, ordinary integers and factorials in
the explicit formulas are replaced by their quantum analogues, for example
\begin{equation}
 [n]_{q_{\rm alg}}
 =\frac{q_{\rm alg}^{n}-q_{\rm alg}^{-n}}
        {q_{\rm alg}-q_{\rm alg}^{-1}},\qquad
 [n]_{q_{\rm alg}}\longrightarrow n\quad(q_{\rm alg}\to1).
 \label{eq:introqinteger}
\end{equation}
Thus $[2]_{q_{\rm alg}}=q_{\rm alg}+q_{\rm alg}^{-1}$ tends to $2$,
and the quantum dimension $d_j=[2j+1]_{q_{\rm alg}}$ tends to
the ordinary dimension in this convention.
The full definition also specifies tensor products and braidings;
a braiding is the prescribed operator for exchanging two factors.
It is not merely a rule for replacing integers. At the usual compact
root $q_{\rm alg}=e^{\pi i/r}$, the finite fusion theory retains the
labels $0,\frac12,\ldots,\frac{r-2}{2}$. A triple must satisfy the
triangle inequalities, integral total spin, and
$a+b+c\le r-2$. Taylor--Woodward give the representation-theoretic
definition and these conventions explicitly~\cite{TaylorWoodward}.

\paragraph{Why these coefficients can build an invariant.}
With four tensor factors there are five ways to put in parentheses.
Each elementary change of parentheses has an $F$-matrix.
Following either route between two parenthesizations must give
the same total change of basis. This equality of matrix products
is the \emph{pentagon identity}. In tetrahedral language it supplies
the identity for replacing two tetrahedra by three inside a
triangulated ball, keeping its boundary fixed. The internal labels
are summed because the intermediate basis is not part of the
boundary data. Together with orthogonality and normalization
identities, this makes appropriately defined state sums independent
of the triangulation.

Schematically, a Turaev--Viro sum for a closed triangulated
three-manifold has the form
\begin{equation}
 Z_r(M)=C_r(T)\sum_{\text{admissible edge labels }\mathbf j}
       \prod_{e\in E(T)}d_{j_e}\,
       \prod_{\Delta\in T}\{6j\}_{\Delta,\mathbf j}.
 \label{eq:introstatesum}
\end{equation}
Here $M$ is the manifold, $T$ its triangulation, and $E(T)$ the
set of edges. A \emph{state} $\mathbf j$ assigns an allowed label
$j_e$ to every edge so that the three labels on every triangular
face are admissible. For one state, multiply a quantum dimension
$d_{j_e}$ for each edge and a $6j$-weight for each tetrahedron
$\Delta$, using its six edge labels. Then sum over all states and
multiply by the prescribed normalization $C_r(T)$.
At the finite root there are only finitely many states.
This formula explains the structure; the exact weights and
normalization belong to the chosen theory~\cite[Sections 2.2--2.3]{ChenYang}.

We can now see the route back to geometry. The representation
calculus first supplies exact rules for assembling local weights
into a global invariant. One can then study what happens as the
labels and quantum parameter vary: particular asymptotic regimes
relate the tetrahedral weights to geometric actions and volumes.
The finite example above explains the local algebra needed for
that question; it does not itself assert a volume formula.
Nor is the state sum already an integral over hyperbolic metrics.
Such an interpretation, including its geometric measure, would
require a further result.

\subsection{The invariants and their volume limits}
\label{sec:introinvariants}

The main objects used below can be distinguished by their inputs and
outputs.
\begin{center}
\small
\begin{tabular}{@{}>{\raggedright\arraybackslash}p{0.23\textwidth}>{\raggedright\arraybackslash}p{0.32\textwidth}>{\raggedright\arraybackslash}p{0.39\textwidth}@{}}
\toprule
Object & Input & Output and construction\\
\midrule
Colored Jones $J_N$ & Knot, color $N$, variable $q_J$
& A Laurent polynomial, from braid or skein calculus.\\[4pt]
Kashaev invariant & Knot and integer $N$
& The specialization $J_N(K;e^{2\pi i/N})$ in the reduced convention.\\[4pt]
Reshetikhin--Turaev (RT) & Closed oriented three-manifold and root/level
& A scalar, from weighted colored-link evaluations of a surgery presentation.\\[4pt]
Turaev--Viro (TV) & Closed three-manifold and root/level
& A scalar state sum, built from tetrahedral $6j$-weights.\\[4pt]
Continuous $b$-$6j$ & Six continuous representation parameters and $b$
& A change-of-basis kernel for a noncompact representation theory.\\[4pt]
Liouville correlation & Marked conformal surface, insertion parameters, and $b$
& A function of the marked points, obtained by averaging field insertions.\\
\bottomrule
\end{tabular}
\end{center}
The first four objects are topological invariants with their prescribed
normalizations. The last two supply analytic kernels and functions
that enter the comparisons in this monograph. In particular, a local
$6j$-weight is not already a three-manifold invariant.

The knot construction begins with a braid whose closure is the knot.
A crossing acts by a braiding matrix on the tensor product of the
strand representations; the Yang--Baxter equation gives the braid
relation. A quantum trace, adjusted to the representation theory,
closes the strands, and the prescribed normalization handles the
moves that change the braid presentation~\cite[Section 2]{MurakamiIntro}.

For the colored Jones polynomial, the color $N$ denotes the
$N$-dimensional representation at generic parameter. We use the
\emph{reduced} normalization $J_N(\text{unknot};q_J)=1$, and specialize
only after forming that polynomial. Murakami--Murakami identify this
specialization with Kashaev's invariant~\cite{MurakamiMurakami}.
For a hyperbolic knot, the volume conjecture is
\begin{equation}
 \lim_{N\to\infty}\frac{2\pi}{N}
 \log|J_N(K;e^{2\pi i/N})|=\Vol(S^3\setminus K).
 \label{eq:introknotvolume}
\end{equation}
The right side is the ordinary three-dimensional hyperbolic volume,
not a volume of moduli space. The assertion concerns a sequence in
which both the color and the root change.

The closed-manifold version used in Section~\ref{sec:dehnresponse}
is the Chen--Yang asymptotic at their root $\zeta_r=e^{2\pi i/r}$,
with $r$ odd:
\begin{equation}
 \frac{4\pi}{r}\log|\mathrm{RT}_r(M;\zeta_r)|\longrightarrow\Vol(M),
 \qquad
 \frac{2\pi}{r}\log\mathrm{TV}_r(M;\zeta_r)\longrightarrow\Vol(M).
 \label{eq:introclosedvolume}
\end{equation}
These statements are conjectural in general, with theorems for
specified families; the hyperbolic fillings of the figure-eight used
later are covered by Wong--Yang~\cite{WongYangDehn}. The doubling
relation between matching RT and TV theories explains the factor
two between their leading logarithms. Exact identities require
consistent choices: generally $\mathrm{TV}_{\mathcal C}$ is the RT
invariant of the categorical center, and the familiar modulus-square
formula holds in the matching unitary modular setting
\cite{TuraevVirelizier}. One should not equate differently normalized
$\mathrm{SU}(2)$ and $\mathrm{SO}(3)$ versions without conversion.

There are two root-of-unity subtleties. First, authors use parameters
differing by a square: in a common convention $q_J=q_{\rm alg}^2$.
Within Chen--Yang's conventions, the volume root $e^{2\pi i/r}$
differs from the standard compact root $e^{\pi i/r}$; the latter
belongs to a different asymptotic problem~\cite{ChenYang}.
Second, even where the numerical roots agree after changing notation,
the colors may not belong to the same theory. At
$q_{\rm alg}=e^{\pi i/N}$, the $N$-dimensional color has
$[N]_{q_{\rm alg}}=0$ and lies outside the retained simple labels
of the level-$N-2$ compact fusion theory. Reducing $J_N$ first and
then specializing is essential. Thus the Kashaev limit should not
be read as an allowed fixed-color observable in that finite unitary
theory.

\subsection{A complete example: the figure-eight knot}
\label{sec:introfigureeight}

For the figure-eight knot $4_1$, the root specialization is an
especially simple positive sum:
\begin{equation}
 \boxed{\quad
 J_N(4_1;e^{2\pi i/N})=\sum_{k=0}^{N-1}P_{N,k},\qquad
 P_{N,k}=\prod_{j=1}^{k}4\sin^2\frac{\pi j}{N},\quad P_{N,0}=1.
 \quad}
 \label{eq:introfigureproduct}
\end{equation}
Kashaev gives the product formula~\cite[eqs. (2.2), (2.5)]{KashaevVolume}.
The following elementary asymptotic proof is presented by
Murakami and credited there to Ekholm~\cite[Section 3.2]{MurakamiIntro}.

The ratio of consecutive products is
\[
 \frac{P_{N,k}}{P_{N,k-1}}=4\sin^2\frac{\pi k}{N}.
\]
It is less than one before $N/6$, greater than one between $N/6$
and $5N/6$, and less than one afterward. The global maximum is
therefore attained at $m_N=\lfloor5N/6\rfloor$, allowing an adjacent
tie. To compare with the initial value $1$, use the elementary
identity $P_{N,N-1}=N^2$: evaluating
$(z^N-1)/(z-1)=\prod_{j=1}^{N-1}(z-e^{2\pi i j/N})$ at $z=1$
and taking absolute squares proves it. Positivity now gives
\[
 P_{N,m_N}\le J_N(4_1;e^{2\pi i/N})\le N P_{N,m_N}.
\]
The two logarithmic bounds differ by $\log N$, which disappears
after division by $N$. The remaining Riemann sum has an integrable
logarithmic singularity at zero, and yields
\begin{equation}
 \lim_{N\to\infty}\frac1N\log J_N(4_1;e^{2\pi i/N})
 =\int_0^{5/6}\log(4\sin^2\pi x)\,\dd x.
 \label{eq:introfigureintegral}
\end{equation}
The endpoint causes no difficulty: subtract $2\log x$ near zero;
the remaining function is continuous and the logarithmic part is
controlled by Stirling's formula.

Define the Lobachevsky function by
\[
 \Lambda(\theta)=-\int_0^\theta\log|2\sin t|\,\dd t.
\]
Its elementary symmetry and duplication identities give
\begin{align}
 \lim_{N\to\infty}\frac{2\pi}{N}
       \log J_N(4_1;e^{2\pi i/N})
 &=-4\Lambda(5\pi/6)=4\Lambda(\pi/6)\notag\\
 &=6\Lambda(\pi/3)=2v_3
  =\Vol(S^3\setminus4_1),
 \label{eq:introfigurevolume}
\end{align}
where $v_3=3\Lambda(\pi/3)$ is the volume of a regular ideal
tetrahedron. The complete figure-eight complement consists of two
such tetrahedra. Its volume is approximately $2.02988321$.

This calculation already displays the mechanism: an exact quantum
sum, a logarithmic classical phase, a dominant critical point,
and a geometric critical value. It also shows why selecting the
critical point matters. For
$f(t)=\int_0^t\log(4\sin^2\pi x)\,\dd x$, the stationary point
$t=1/6$ is a minimum, while $t=5/6$ is the relevant maximum.
Here positivity controls the entire sum. Complex-valued sums and
integrals require a separate analysis of their contributing
critical points and possible cancellation.

\subsection{Why dilogarithms occur}
\label{sec:introdilog}

An elementary convergent product explains the terminology. For
$\varepsilon>0$ and $|z|<1$, set
\[
 (z;e^{-\varepsilon})_\infty
      =\prod_{n=0}^\infty(1-ze^{-n\varepsilon}).
\]
Expanding the logarithm and summing the geometric progression gives
\begin{align}
 \log(z;e^{-\varepsilon})_\infty
 &=-\sum_{m=1}^\infty
       \frac{z^m}{m(1-e^{-m\varepsilon})}\notag\\
 &=-\frac{\operatorname{Li}_2(z)}{\varepsilon}
       +\frac12\log(1-z)+O(\varepsilon),
 \label{eq:introqpochhammer}
\end{align}
uniformly on compact subsets of $|z|<1$, with the logarithm
continued from $z=0$. Indeed,
$(1-e^{-t})^{-1}-t^{-1}-\tfrac12=O(t)$ for all $t>0$;
for $|z|\le\rho<1$ the remainder is bounded by a constant times
$\varepsilon\sum_{m\ge1}\rho^m$. This justifies summing the
expansion despite unbounded $m$.
The logarithm of this ``quantum dilogarithm'' has the ordinary
dilogarithm as its leading term. Faddeev--Kashaev's quantum
pentagon identity deforms the classical five-term dilogarithm
identity~\cite{FaddeevKashaev}.

The geometry enters in two ways. For an oriented ideal tetrahedron
with shape $z$ in the upper half-plane,
\begin{equation}
 \Vol(T_z)=D(z)
 =\Im\operatorname{Li}_2(z)+\arg(1-z)\log|z|.
 \label{eq:introBlochWigner}
\end{equation}
The logarithmic correction matters: an arbitrary dilogarithm phase
does not have volume as its imaginary part without the appropriate
boundary terms. Also,
$\dd\operatorname{Li}_2(z)=-\log(1-z)\,\dd\log z$.
Consequently differentiation of a sum of dilogarithms produces
logarithms of tetrahedral shapes, which can impose gluing equations.

Here is a particularly direct version in the convention of
Dimofte--Garoufalidis~\cite[Section 5.1]{DimofteGaroufalidis}.
On a fixed logarithmic branch away from the limiting pole cuts,
their tetrahedral quantum function obeys
\begin{equation}
 \psi_\hbar(Z+\hbar)=(1-e^{-Z})\psi_\hbar(Z),\qquad
 \log\psi_\hbar(Z)\sim\frac1\hbar\operatorname{Li}_2(e^{-Z})+O(1).
 \label{eq:introquantumshape}
\end{equation}
Putting $\psi_\hbar\sim e^{S/\hbar}$ in the difference equation
gives $e^{S'(Z)}=1-e^{-Z}$. With $z=e^Z$, the right side is
the companion tetrahedral shape $1-z^{-1}$. Thus a difference
equation for a special function becomes an algebraic shape relation
in the limit.

A \emph{state integral} combines these tetrahedral functions with
quadratic exponential factors determined by the gluing data and
integrates over internal logarithmic variables. Differentiating its
leading phase gives the logarithmic gluing equations; exponentiating
them gives the familiar products of shapes around edges. The
boundary variables are retained as arguments of the integral.
The precise quadratic term, logarithm branches, and contour are
part of the construction. The convergent product in
\eqref{eq:introqpochhammer} explains the mechanism in a simple
domain; one cannot substitute a unit-modulus root into that product
and assume it converges. The continued or noncompact functions used
later have their own definitions and contour prescriptions.
\subsection{Stationary phase, loops, torsion, and degenerate saddles}
\label{sec:introsemiclassical}

The word \emph{semiclassical} means an asymptotic regime in which
the quantum parameter tends to zero while the intended geometric
data remain controlled. Much of the calculation is ordinary
asymptotic analysis. Consider
\begin{equation}
 I(\hbar)=\int_\Gamma e^{S(x)/\hbar}
                 a(x;\hbar)\,\dd^m x,
 \qquad a(x;\hbar)=a_0(x)+\hbar a_1(x)+\cdots.
 \label{eq:introsaddleintegral}
\end{equation}
Suppose $S$ is holomorphic near a nondegenerate critical point
$x_*$, $a_0(x_*)\ne0$, and the amplitude has a classical-symbol
expansion there, including the derivatives needed at each order.
Let $\hbar$ tend to zero along a fixed ray, and localize to an
oriented steepest-descent contour through $x_*$ on which the real
part of $(S-S(x_*))/\hbar$ has a strict quadratic maximum.
Taylor expansion, rescaling, and Gaussian integration give this
localized contribution
\begin{equation}
 I_{x_*}(\hbar)\sim
 e^{S(x_*)/\hbar}(2\pi\hbar)^{m/2}
 \frac{a_0(x_*)}{\sqrt{\det(-S''(x_*))}}
 (1+c_1\hbar+\cdots).
 \label{eq:introgaussian}
\end{equation}
The contour orientation, the ray of $\hbar$, and the continued
square root fix the phase. Extending the formula from this local
integral to the whole contour requires a bound on its complement.
For a real
Laplace integral with a strict maximum this is the usual Gaussian
formula; for a complex phase it is applied along suitable real
$m$-dimensional contours in the complexified variables.

This formula gives a useful translation of the physics terminology.
The \emph{classical} or \emph{tree-level} contribution is the
critical value $S(x_*)/\hbar$. The \emph{one-loop} contribution
includes the Gaussian determinant and leading amplitude or measure
factors. Higher-loop terms come from the higher Taylor derivatives
and amplitude corrections. Feynman diagrams organize the Gaussian
contractions needed to compute these coefficients; the terminology
does not require an additional physical assumption. In logarithmic
form one writes
\[
 \log I_{x_*}=\frac{S_0}{\hbar}
      +\nu\log\hbar+S_1+S_2\hbar+\cdots,
\]
where $\nu$ depends on the integral's normalization and any zero
modes (kernel directions of the linearized operator). A geometric expansion of $S_0$ in a shrinking area or an
inverse Dehn-filling slope is a separate expansion. The cubic and
quintic Neumann--Zagier coefficients are not, on that account,
quantum higher-loop coefficients.

\emph{Torsion} is the topological form taken by a determinant of
linearized equations after gauge directions have been accounted
for. For a flat holonomy representation $\rho$, the relevant
cellular complex is $C^\bullet(M;\operatorname{Ad}\rho)$, with the
holonomy acting on the Lie algebra by conjugation.
Its $H^0$ records infinitesimal stabilizers and its $H^1$ records
infinitesimal deformations. Reidemeister torsion compares the
specified cellular bases with bases adapted to the differentials,
using an alternating product of determinants. Nonzero cohomology
requires cohomology bases as well; for a cusped three-manifold,
peripheral curves enter this normalization. This explains why adjoint torsion is a candidate for a
semiclassical amplitude~\cite{WittenJones,DimofteGaroufalidis}.
In particular, the identification of the gluing-equation one-loop
invariant with adjoint Reidemeister torsion is a conjecture in
\cite[Conjecture~1.8]{DimofteGaroufalidis}; its formal Gaussian
construction is not a general proof of that identification.

An individual Hessian determinant is not the torsion by definition:
measure factors, gauge directions, and basis conventions must be
matched. In the surgery convention used below the relevant
amplitude is proportional to $\operatorname{Tor}^{-1/2}$.
Section~\ref{sec:dehnresponse} supplies the actual peripheral and
solid-torus factors and an established figure-eight comparison.

There is also a global question. A \emph{Lefschetz thimble} is a
steepest-descent cycle attached to a critical point. Under the
usual analytic hypotheses, the exact contour can be represented
in an appropriate relative homology group as
\[
 \Gamma=\sum_\alpha n_\alpha\mathcal J_\alpha.
\]
The coefficients determine which saddles contribute. Their
description can change across a \emph{Stokes locus} as parameters
vary; the exact integral can remain analytic while its decomposition
into saddle contributions changes~\cite[Sections 2--3]{WittenContinuation}.
To recover the asymptotic of a whole invariant, one must control
these coefficients, all competing contributions, and cancellation.
The figure-eight positive sum supplied this global control by an
elementary inequality. A local geometric saddle in a general
state integral supplies only its local contribution until these
additional questions are settled.

Finally, a degenerate Hessian changes the local model. At a simple
fold, two saddles merge. After separating the nondegenerate
transverse directions, the remaining phase has the form
\[
 S(x;t)=S_0+\frac{x^3}{3}-tx+\text{higher terms}.
\]
The substitution $x=\hbar^{1/3}y$, $t=\hbar^{2/3}s$ gives
\[
 \frac{S-S_0}{\hbar}=\frac{y^3}{3}-sy+\cdots.
\]
The resulting contour integral is an Airy integral, rather than
a sum of two separately uniform Gaussian approximations. This
explains the exponent $2/3$ in the transition scale. For the regular
family in Section~\ref{sec:quantumAiry}, the real small parameter
is $h=2\pi b^2$, so the geometric deformation has scale
$b^{4/3}$. When that deformation is proportional to boundary area
$A$, the same scale follows from $W\asymp A^{3/2}$ and $W/h=O(1)$.
The later theorem proves the contour and additive error estimates
for its pole-cleared regular family and fixed bounded real
transition parameter. Its volume-action parameter is
$\hbar_{\mathrm{vol}}=h/2$ as in~\eqref{eq:TVvolumehbar}; this
constant does not change the exponent $2/3$. The scaling argument
alone does not prove a transition theorem.
Figure~\ref{fig:mono-airy} illustrates the cubic normal form and its
stationary points; its parameter $t$ is the local fold parameter above.

% BEGIN INLINE FIGURE: 09_airy.tex
\begin{figure}[htbp]
\centering
\begin{tikzpicture}[font=\small,>=Latex,line cap=round]
  \foreach \shift/\mval/\heading in {-4.5/-1.4/{$t<0$},0/0/{$t=0$},4.5/1.4/{$t>0$}} {
    \begin{scope}[shift={(\shift,0)}]
      \node[font=\bfseries\small] at (0,2.0) {\heading};
      \begin{scope}[x=.78cm,y=.5cm]
        \draw[->,mpgray!60] (-2.3,0)--(2.4,0) node[below,font=\scriptsize] {$x$};
        \draw[->,mpgray!60] (0,-3.2)--(0,3.2);
        \begin{scope}
          \clip (-2.2,-3) rectangle (2.2,3);
          \draw[mpblue,line width=1.15pt,domain=-2.2:2.2,samples=150]
            plot (\x,{\x*\x*\x/3-\mval*\x});
        \end{scope}
      \end{scope}
    \end{scope}
  }
  \filldraw[draw=mpochre,fill=white,line width=1.2pt] (0,0) circle (3pt);
  \begin{scope}[shift={(4.5,0)},x=.78cm,y=.5cm]
    \draw[mpgray!55,dashed] (-1.183216,0)--(-1.183216,1.104334)
      (1.183216,0)--(1.183216,-1.104334);
    \fill[mpochre] (-1.183216,1.104334) circle (2.3pt);
    \fill[mpochre] (1.183216,-1.104334) circle (2.3pt);
    \node[below left=3pt,font=\scriptsize] at (-1.183216,0) {$-\sqrt t$};
    \node[below=3pt,font=\scriptsize] at (1.183216,-1.104334) {$+\sqrt t$};
  \end{scope}
  \node[align=center,text=mpgray,font=\scriptsize] at (-4.5,-2.05) {No real stationary point};
  \node[align=center,text=mpgray,font=\scriptsize] at (0,-2.05) {One double stationary point};
  \node[align=center,text=mpgray,font=\scriptsize] at (4.5,-2.05) {Two simple stationary points};
  \node[align=center] at (0,-2.85)
    {$\mathcal S_t(x)=\displaystyle\frac{x^3}{3}-t x,
      \qquad x=O(\hbar^{1/3}),\quad t=O(\hbar^{2/3})$.};
\end{tikzpicture}
\caption[Coalescing stationary points and the Airy normal form]{The cubic normal form for an Airy transition. The stationary
equation is $x^2=t$; two stationary points merge at $t=0$.
The scaling shown balances the cubic and linear terms in
$\exp(i\mathcal S_t/\hbar)$. This is a normal-form illustration,
not the full phase or integration contour of a quantum invariant;
normalization, amplitudes, and contour tails require the separate
estimates in the text.}
\label{fig:mono-airy}
\end{figure}
% END INLINE FIGURE

\subsection{Flat connections, Chern--Simons theory, and boundary states}
\label{sec:introCS}

The connection with hyperbolic geometry can be phrased using flat
connections. If $\mathscr A$ is a connection one-form, its curvature
is $F_{\mathscr A}=\dd\mathscr A+\mathscr A\wedge\mathscr A$.
In a chosen trivialization the Chern--Simons functional has the form
\begin{equation}
 \operatorname{CS}(\mathscr A)=\frac1{4\pi}\int_M
   \tr\left(\mathscr A\wedge\dd\mathscr A
                    +\frac23\mathscr A\wedge\mathscr A\wedge\mathscr A\right).
 \label{eq:introCSfunctional}
\end{equation}
The trace convention fixes the periods and overall normalization.
On a closed manifold its first variation is
$\delta\operatorname{CS}=(2\pi)^{-1}\int_M
\tr(\delta\mathscr A\wedge F_{\mathscr A})$.
Thus the critical points are flat connections modulo gauge.
The formal expression
$\int e^{ik\operatorname{CS}(\mathscr A)}\,D\mathscr A$
motivates the theory; exact algebraic constructions of quantum
invariants do not depend on treating this symbol as an ordinary
convergent integral~\cite{WittenJones}.

For the compact group $\mathrm{SU}(2)$, the flat connections form
the compact character-variety problem. Hyperbolic three-dimensional
holonomy instead takes values in $\mathrm{PSL}_2(\CC)$, with a lift
to $\mathrm{SL}_2(\CC)$ when chosen. It supplies a distinguished
\emph{complex} flat connection. The correctly normalized complex
Chern--Simons action contains both the volume and the Chern--Simons
invariant. Analytic continuation and the choice of integration
cycle connect these settings; merely changing the name of the
gauge group does not specify the cycle~\cite{GukovApoly,WittenContinuation}.

For a manifold with boundary, variation of the functional has a
boundary term. This is the source of the symplectic form on boundary
flat connections and the need to specify boundary data. Abstractly,
a topological quantum field theory assigns a vector space
$\mathcal H(\Sigma)$ to an oriented surface, and a vector
$\Psi_M\in\mathcal H(\partial M)$ to a three-manifold with boundary.
Gluing is a pairing:
\[
 Z(M_1\cup_\Sigma M_2)=\langle\Psi_{M_1},\Psi_{M_2}\rangle,
\]
with orientation reversal supplying the dual space and the theory's
framing conventions understood. In finite root-of-unity theories
this is finite-dimensional linear algebra. In the analytic models
used for complex holonomy, a state is represented locally by a
function or distribution of continuous boundary variables, and
the pairing becomes an integral. Those analytic models require
their own existence and convergence statements.

A \emph{polarization} chooses which half of the conjugate boundary
coordinates are arguments of the state. For a torus boundary,
use the multiplier logarithms $u,v$ of Section~\ref{sec:dehnresponse}.
The Neumann--Zagier normalization is
\begin{equation}
 \Omega_\partial=\frac12\dd v\wedge\dd u,
 \qquad \alpha_\partial=\frac12v\,\dd u.
 \label{eq:introNZpolarization}
\end{equation}
The geometric deformations of the interior give a Lagrangian curve
$v=v(u)$. A semiclassical boundary state associated with this
branch has the form
\begin{equation}
 \Psi_M(u;\hbar)\sim e^{\Phi(u)/\hbar}
       (a_0(u)+\hbar a_1(u)+\cdots),\qquad
 \Phi'(u)=\frac12v(u).
 \label{eq:introWKBstate}
\end{equation}
The leading exponent is a primitive of the symplectic potential
on the Lagrangian curve. This is the geometric content of the WKB
(Wentzel--Kramers--Brillouin) ansatz~\cite{DimofteQuantumRS}.

This choice of half-boundary variables also explains the boundary
correction to volume. With $F'=v$ and $F(0)=0$, the note uses
$\Phi=\Phi_0+F/2$, with $\Im\Phi_0=V_0$ the complete volume, while
\begin{equation}
 V(u)-V_0=\frac12\Im F(u)-\frac12\Re u\,\Im v(u).
 \label{eq:introvolumeboundary}
\end{equation}
The last term is essential. The imaginary part of an action in an
arbitrary boundary polarization need not itself be the incomplete
hyperbolic volume. Section~\ref{sec:NZsymplectic} derives this
identity and the corresponding filling transformation explicitly.

\subsection{Quantizing peripheral holonomy and performing surgery}
\label{sec:introquantization}

Set $P=v/2$, so that $\Omega_\partial=\dd P\wedge\dd u$.
In the formal convention compatible with
$e^{\Phi/\hbar}$, use operators
\begin{equation}
 \widehat u\,\Psi=u\Psi,\qquad
 \widehat P\,\Psi=\hbar\partial_u\Psi,
 \qquad [\widehat P,\widehat u]=\hbar.
 \label{eq:introcommutator}
\end{equation}
The usual physical factor $i$ has been absorbed into the parameter
$\hbar$; the formula defines the convention here. Acting on the
leading exponential gives $\widehat P\Psi\sim\Phi'\Psi=P\Psi$.
Thus differentiation becomes the classical longitude coordinate
in the small-$\hbar$ limit.

In local eigenvalue coordinates $m=e^{u/2}$ and $l=e^{v/2}$,
the corresponding operators are multiplication and translation:
\[
 \widehat m\Psi(u)=e^{u/2}\Psi(u),\qquad
 \widehat l\Psi(u)=\Psi(u+\hbar),\qquad
 \widehat l\widehat m=e^{\hbar/2}\widehat m\widehat l.
\]
Fixed signs from an $\mathrm{SL}_2$ lift can multiply $m$ and $l$;
they do not change this commutation relation. The factor
$e^{\hbar/2}$ is the Weyl parameter in these coordinates. It should
not be identified with every author's variable called $q$.

The classical boundary holonomies satisfy a polynomial relation
$A(l,m)=0$ on the appropriate character-variety component. To
\emph{quantize} such a relation is to seek a difference operator
$\widehat A(\widehat l,\widehat m;\hbar)$ annihilating the boundary
state. Substituting \eqref{eq:introWKBstate} gives, to leading order,
\[
 A\bigl(e^{\Phi'(u)},e^{u/2}\bigr)=0.
\]
This recovers the classical curve because $\Phi'=v/2$. It does not
uniquely determine the exact operator: ordering and lower-order
terms remain. Identifying a particular operator with a recursion
for a particular quantum invariant is a further mathematical
statement~\cite{GukovApoly,DimofteQuantumRS}.

Dehn filling has a similarly concrete meaning. For a slope $(p,q)$,
choose $\rho,\sigma\in\ZZ$ with $p\sigma-q\rho=1$ and put
\[
 a=pu+qv,\qquad b_{\rm core}=\rho u+\sigma v.
\]
Then $\dd b_{\rm core}\wedge\dd a=\dd v\wedge\dd u$. For $q\ne0$
the generating function
\begin{equation}
 G(u,a)=\frac{pu^2}{4q}-\frac{ua}{2q}
                    +\frac{\sigma a^2}{4q}
 \label{eq:introfillingkernel}
\end{equation}
satisfies $\dd G=\frac12b_{\rm core}\,\dd a-\frac12v\,\dd u$ on
this relation. An integral transform with exponential kernel
$e^{G/\hbar}$ therefore changes the boundary coordinate from $u$
to $a$. Its stationary equation for $\Phi(u)+G(u,a)$ is exactly
$a=pu+qv$. At the classical level, pairing with a solid-torus state imposes
the geometric filling branch $a=2\pi i$; the cone-angle family
uses $a=i\theta$.
The full surgery kernel includes amplitude and topological
normalization factors in addition to this exponential.

In the closed-filling convention of this monograph,
$\hbar_r=4\pi i/r$. If the filled critical phase has imaginary
part $V$, then
\begin{equation}
 \left|e^{\Phi_{\rm filled}/\hbar_r}\right|
       =\exp\left(\frac{rV}{4\pi}\right).
 \label{eq:introRTgrowth}
\end{equation}
This is the RT exponential in \eqref{eq:introclosedvolume}.
After the Fourier-kernel square-root normalization, the Gaussian
change of boundary coordinate has dependence
$(p+qv'(u))^{-1/2}$, up to its phase; the solid torus contributes
further factors. This requires $q\ne0$ in the displayed chart and
$p+qv'(u)\ne0$ at the saddle.
The calculation in Section~\ref{sec:dehnresponse} makes each of
these statements precise and states the hypotheses needed to pass
from the geometric contribution to an exact invariant. Notice that
the inverse level $r^{-1}$ and the inverse slope length
$|p+q\tau|^{-1}$ are independent small parameters; here
$\tau=v^\prime(0)$ is the complete cusp shape.

\subsection{Liouville theory, conformal blocks, and continuous fusion}
\label{sec:introLiouville}

Take $0<b<1$. The other quantum construction in this monograph starts in two
dimensions. In Liouville theory the variable is a conformal factor
$\phi$ on a fixed Riemann surface with reference metric $g$;
formally the new area element is $e^{2b\phi}\dd A_g$.
The field is distributional, so this exponential and its point
insertions require renormalization. The Gaussian free field is a
generalized Gaussian function with inverse-Laplacian covariance.
Gaussian multiplicative chaos constructs its exponential measure
by smoothing the field, subtracting half the variance of the
scaled field from the exponent, and taking a limit. In the subcritical range above this
gives the rigorous random measures used here~\cite{KRV}.

The basic observables are the renormalized exponential insertions
$V_\alpha(p)\sim e^{2\alpha\phi(p)}$. A \emph{correlation function}
is the averaged product
\[
 C_b(\boldsymbol\alpha;\boldsymbol p)
       =\left\langle\prod_i V_{\alpha_i}(p_i)\right\rangle.
\]
It depends on the marked conformal surface. It is neither a
three-manifold invariant nor, without a further operation, a
measure on moduli space. The parameters governing conformal
symmetry are
\begin{equation}
 Q=b+b^{-1},\qquad c_L=1+6Q^2,\qquad
 \Delta(\alpha)=\alpha(Q-\alpha).
 \label{eq:introLiouvilleparameters}
\end{equation}
Here $c_L$ is the \emph{central charge} and $\Delta$ the
\emph{conformal weight}. Concretely the infinitesimal symmetry
algebra is the Virasoro algebra
\[
 [L_m,L_n]=(m-n)L_{m+n}
       +\frac{c_L}{12}(m^3-m)\delta_{m+n,0}.
\]
A highest-weight vector has $L_0$-eigenvalue $\Delta$ and is killed
by $L_n$ for $n>0$; its descendants are obtained with the negative
generators. This representation theory determines the power
series called conformal blocks~\cite{TeschnerRevisited}.

On a four-marked sphere with marks $0,1,z,\infty$, a pants
decomposition first combines two of the insertions. Fixing the
intermediate representation with label $\alpha_s$ gives a
\emph{chiral conformal block}: a holomorphic, generally
multivalued function with a specified leading behavior, for example
\[
 \mathcal F^{(s)}_{\alpha_s}(z)
       =z^{\Delta(\alpha_s)-\Delta(\alpha_1)-\Delta(\alpha_3)}
                       (1+O(z)).
\]
``Chiral'' refers here to the holomorphic half. The full
correlation pairs holomorphic and antiholomorphic blocks and
integrates over the intermediate spectrum. Label the insertions at $0,1,z,\infty$ by
$\alpha_1,\alpha_2,\alpha_3,\alpha_4$, respectively.
For real external charges in the initial domain
$\alpha_i<Q/2$, $\alpha_1+\alpha_3>Q/2$, and
$\alpha_2+\alpha_4>Q/2$, the proved four-point decomposition has
the form
\[
 C_b(z,\bar z)=\int_0^\infty
    \rho(\alpha_1,\ldots,\alpha_4;P)\,
    |\mathcal F^{(s)}_{Q/2+iP}(z)|^2\,\dd P,
\]
with normalization and products of three-point structure constants
included in $\rho$. The explicit three-point constants are the
DOZZ constants, named for Dorn--Otto and Zamolodchikov--Zamolodchikov.
The relevant three-point formula and four-point decomposition have
rigorous constructions and proofs in their stated domains
\cite{KRV,GKRVBootstrap}. Here the charges in~\cite[Theorem~1.1]{GKRVBootstrap} are twice
our $\alpha_i$. Continuing outside that domain can require
additional residues when poles cross the spectral contour.
For complex charges the two chiral factors must be continued
analytically in the charges; the real-domain absolute-square
notation does not prescribe that continuation.

Changing the pants decomposition gives the \emph{fusion kernel}:
\begin{equation}
 \mathcal F^{(s)}_{\alpha_s}(z)
   =\int_{\mathcal C} F_{\alpha_s\alpha_t}\,
             \mathcal F^{(t)}_{\alpha_t}(z)\,\dd\alpha_t.
 \label{eq:introfusion}
\end{equation}
The measure is either displayed separately or incorporated in $F$.
This is the continuous counterpart of
\eqref{eq:introrecoupling}: four external labels and two internal
ones again give six labels. Rescaling the basis blocks changes
the displayed kernel. Accordingly the fusion kernel is related
to a $b$-$6j$ symbol through explicit trivalent normalization and
spectral factors, not by an unqualified equality of scalar
functions~\cite[Section 4]{TeschnerVartanov}.

The noncompact representation theory behind this kernel is the
\emph{modular double}. It uses the compatible deformation parameters
\[
 q_{\rm alg}=e^{\pi i b^2},\qquad
 \widetilde q_{\rm alg}=e^{\pi i/b^2}.
\]
The principal representation labels are $\alpha=Q/2+iP$; tensor
products decompose by direct integrals instead of finite sums.
The double refers to the simultaneous roles of $b$ and $b^{-1}$.
The corresponding $b$-$6j$ is a continuous recoupling kernel.
This is a different setting from the finite root-of-unity label
set in Section~\ref{sec:introrecoupling}.

Its explicit formulas use Barnes' double Gamma function $\Gamma_b$,
a meromorphic analogue of the Gamma function with the two step
sizes $b$ and $b^{-1}$, and the \emph{double sine}
\begin{equation}
 S_b(x)=\frac{\Gamma_b(x)}{\Gamma_b(Q-x)},\qquad
 \begin{aligned}
 S_b(x+b)&=2\sin(\pi bx)S_b(x),\\
 S_b(x+b^{-1})&=2\sin(\pi x/b)S_b(x),\\
 S_b(x)S_b(Q-x)&=1.
 \end{aligned}
 \label{eq:introdoublesine}
\end{equation}
These identities explain why sine factors and pole movements occur
in the later calculations. The double sine is related to Faddeev's
noncompact quantum dilogarithm by shifted arguments and explicit
Gaussian factors; conventions matter~\cite[Appendix A]{TeschnerVartanov}.

The first shift relation also explains its classical limit. Put
$L_b(x)=b^2\log S_b(x/b)$ for real $x$ in the fundamental strip.
Then
\[
 L_b(x+b^2)-L_b(x)=b^2\log(2\sin\pi x).
\]
A smooth leading limit must therefore have derivative
$\log(2\sin\pi x)$. Reflection fixes the constant, giving
$L(x)=-\Lambda(\pi x)/\pi$ for $0<x<1$. This is a formal
derivation of the leading term; Lemma~\ref{lem:stripdoubleSine}
supplies the uniform estimate used in the proof. Thus the same
Lobachevsky function that appeared in the elementary figure-eight
sum also appears in the continuous $b$-$6j$ analysis.

\subsection{Why a screening residue meets flat cone geometry}
\label{sec:introscreening}

The particular Liouville specialization in this monograph can be
understood through a one-variable integral. Separate the constant
part $c_0$ of the field from its nonconstant part. After the
insertions have shifted the Gaussian field, let $Z>0$ be the
remaining random area factor and put
\[
 s=\frac{\sum_i\alpha_i-Q}{b}.
\]
With the overall normalization adopted in
\eqref{eq:LiouvilleGMC}, the integral over the constant mode is
\begin{equation}
 2\int_{\RR}e^{2bs c_0-\Lambda_L Z e^{2bc_0}}\,\dd c_0
 =\frac1b\Lambda_L^{-s}\Gamma(s)Z^{-s},\qquad \Re s>0.
 \label{eq:introzeromode}
\end{equation}
This is just the substitution $t=\Lambda_L Z e^{2bc_0}$. The
positive number $\Lambda_L$ is the cosmological constant, or
coefficient penalizing total area. The formula explains both the
Gamma factor in the correlation and the usefulness of removing
that factor in the unit-volume normalization.

Now take the deficit weights of the flat cone metric:
\[
 0<\mu_i<1,\quad \sum_i\mu_i=2,\qquad
 \alpha_i=\frac{\mu_i}{2b}.
\]
Then $\sum_i\alpha_i=1/b=Q-b$, and therefore $s=-1$.
The raw integral \eqref{eq:introzeromode} does not converge there.
Its meromorphic continuation has a Gamma pole, whereas the
continued unit-volume expression contains $Z^{-s}=Z$: its
\emph{first moment}. The normalized random measure has the prescribed
reference measure as its expectation. The reference factors cancel
at these weights, leaving, in a coordinate with all marks finite,
\begin{equation}
 \mathbb E[Z]=\mathcal A_\mu(\boldsymbol p)
   =\int_{\CC}\prod_i|w-p_i|^{-2\mu_i}\,\dd^2w.
 \label{eq:introfirstmoment}
\end{equation}
This is precisely the area of the flat cone metric that underlies
the Thurston--Deligne--Mostow construction. The charge convention
in the detailed probabilistic formulas is $\beta_i=2\alpha_i$.
Section~\ref{sec:quantummeasure} checks the continuation and
integrability conditions and proves the exact finite-$b$ identity.

The remaining factors in the correlation are pairwise powers of
$|p_i-p_j|$, whose logarithms are pluriharmonic away from collisions.
Consequently the operation that produces the geometric form is
\begin{equation}
 -2i\,\partial\bar\partial\log C_b^{\rm uv}
     =\omega_{\mathrm{CH},\mu},\qquad
 \dd V_{\mathrm{CH},\mu}
     =\frac{\omega_{\mathrm{CH},\mu}^{n-3}}{(n-3)!}.
 \label{eq:introcurvaturemeasure}
\end{equation}
The correlation supplies a logarithmic potential; its logarithmic
curvature supplies the form; the top exterior power supplies the
moduli measure. These are three different operations. This is why
the exact comparison does not identify an unmodified quantum
amplitude with a WP or complex hyperbolic volume density.

The name \emph{one screening} has an equally concrete explanation.
The insertion $V_b$ has conformal weight
$\Delta(b)=b(Q-b)=1$, so integrating it gives a conformally
covariant insertion. At $\sum_i\alpha_i=Q-b$, one such insertion
balances the background charge. Its chiral contour integrals are
\[
 \int_\gamma\prod_i(w-p_i)^{-\mu_i}\,\dd w,
\]
the Deligne--Mostow periods. For four points the relevant twisted
period space is two-dimensional, and changing contour bases gives
the finite hypergeometric fusion matrix calculated in
Section~\ref{sec:screenedfusion}. Generic Liouville fusion is still
the integral operator \eqref{eq:introfusion}; the $2\times2$ matrix
is the exact change of basis for these finite periods. Identifying
it with a residue of the continued infinite fusion operator also
requires the spectral-continuation and tail input stated in
Section~\ref{sec:screenedfusion}. The first-moment area identity
does not require that additional spectral input. At balanced
weights a logarithmic basis is needed. The auxiliary tetrahedron
also need not be the convex realization of the original cone
sphere. Section~\ref{sec:normalizedsixj} checks that geometric
comparison separately.

\subsection{Parameters, limits, and a route through the calculations}
\label{sec:introreadingguide}

The notation records several independent limits. The following
dictionary prevents a formal resemblance of symbols from
identifying different asymptotic questions.
\begin{center}
\small
\begin{tabular}{@{}>{\raggedright\arraybackslash}p{0.23\textwidth}>{\raggedright\arraybackslash}p{0.71\textwidth}@{}}
\toprule
Parameter & Meaning and controlled limit\\
\midrule
$N$ & Color and root order in the Kashaev specialization;
both change in \eqref{eq:introknotvolume}.\\[4pt]
$r$ & Root-order parameter in the finite three-manifold
invariants; $\hbar_r=4\pi i/r$ in the closed-filling phase here.\\[4pt]
$b$ & Continuous Liouville/modular-double parameter;
the regular-family calculation uses $h=2\pi b^2$.\\[4pt]
$\alpha_i$ & Representation or insertion labels; the flat-weight
Liouville specialization uses $\alpha_i=\mu_i/(2b)$.\\[4pt]
$A$ & Geometric boundary area tending to zero in the
Euclidean degeneration; it is not a quantum parameter.\\[4pt]
$(p,q)$ and $\theta$ & Dehn-filling slope and angular holonomy parameter;
$pu+qv=i\theta$. It is the core cone angle when a cone completion exists;
$\theta=2\pi$ gives smooth filling when geometrically realized.\\
\bottomrule
\end{tabular}
\end{center}
The last charge prescription is the pure heavy-charge scaling:
the labels grow like $b^{-1}$ while the deficit weights stay fixed.
The exact length-kernel dictionary has additional finite-$b$
shifts, recorded in \eqref{eq:b6jparameterdictionary}.
Fixed labels with $q\to1$, labels growing with the order of a root
of unity, and the continuous $b\to0$ limit are different problems.
The classical large-spin $6j$ asymptotic involves Euclidean
tetrahedra~\cite{Roberts}; particular compact root regimes involve
spherical tetrahedra~\cite{TaylorWoodward}; the hyperbolic regimes
used below have their own label and contour hypotheses. Geometric
edge lengths or angles are assigned only after specifying the
relevant regime.

The quantum part of this monograph can now be read along three routes.
For the boundary-moduli comparison, begin with
Section~\ref{sec:quantummeasure}, then follow the four-point residue
and its normalization in Sections~\ref{sec:screenedfusion}
and~\ref{sec:normalizedsixj}. For the degeneration of an individual
tetrahedron, use Section~\ref{sec:quantumdegeneration}, the regular
family in Section~\ref{sec:normalizedsixj}, and the uniform theorem for the cleared regular-family integral
at bounded real Airy parameter in Section~\ref{sec:quantumAiry}. For three-manifold volume
and surgery, Section~\ref{sec:dehnresponse} is largely independent
of the Liouville calculations; its symplectic explanation is in
Section~\ref{sec:NZsymplectic}.

An exact identity of kernels, a local saddle expansion, and a
theorem about a global invariant answer different questions. The
later sections state which one is obtained at each step. This
distinction permits the geometric calculations to be used fully
without silently assuming an integration cycle, a measure on
moduli, or a dominance theorem that has not been established.

\paragraph{Computational support.}
The finite companion in \path{experiments/quantum_geodesics/README.md}
connects groups Q1, Q3 and Q4 to the figure-eight product sum,
peripheral/filling conventions and charge/curvature dictionary above.
The result path is \path{artifacts/quantum_geodesics/result.json};
\path{book/claims/sections10-11.json} maps the groups to the precise claims.
Exact finite algebra and selected multiprecision comparisons are separate
records. They do not establish a general colored-Jones or $6j$ asymptotic,
a global quantum surgery contribution, a spectral-tail bound or a GMC
construction. Source identities, implementation review and generation
are recorded separately in the final-chapters checkpoint.
